% ============================================================
% f2014_hunan_science_q15_squares_parabola.tex
% 2014 湖南卷（理科）第 15 题：正方形 ABCD、DEFG 与抛物线 y^2=2px
% 来源 PDF：普通高考/2014/2014湖南理.pdf
% 原页：第 1 页（A3 横排）第 3 栏顶部，印面 885，第 15 题题干上方
% 结构：O 为原点，AD 在 x 轴上且 O 为 AD 中点，|AD|=a=1；
%       ABCD 在 x 轴下方，DEFG 在 x 轴上方且以 D 为左顶点，|DE|=b=1+√2；
%       C、F 在抛物线 y^2=2px（p=a=1）上。比例 1.2 cm/单位。
% ============================================================
\begin{tikzpicture}[>=Stealth, line width=0.6pt, font=\small,
  x=1.2cm, y=1.2cm, line cap=round, line join=round]
  % 坐标轴
  \draw[->] (-1.05,0) -- (3.75,0) node[right] {$x$};
  \draw[->] (0,-2.45) -- (0,2.90) node[above] {$y$};
  % 抛物线 y^2=2x
  \draw[domain=-2.65:2.65, samples=150, smooth, variable=\t]
    plot ({(\t*\t)/2}, {\t});
  % 正方形 ABCD（x 轴下方）
  \draw (-0.5,0) -- (-0.5,-1) -- (0.5,-1) -- (0.5,0) -- cycle;
  % 正方形 DEFG（x 轴上方）
  \draw (0.5,0) -- (2.91421,0) -- (2.91421,2.41421) -- (0.5,2.41421) -- cycle;
  % 原点与顶点标注
  \node[below left]  at (-0.06,-0.04) {$O$};
  \node[below left]  at (-0.5,-0.06) {$A$};
  \node[below left]  at (-0.5,-1.02) {$B$};
  \node[right]       at (0.56,-0.96) {$C$};
  \node[below right] at (0.5,-0.06) {$D$};
  \node[below]       at (2.91421,-0.10) {$E$};
  \node[above right] at (2.88000,2.47000) {$F$};
  \node[above left]  at (0.52000,2.47000) {$G$};
\end{tikzpicture}
